\section{Plug and Process Loads}
\label{sec:plug_and_process_loads}

Plug and process loads (PPLs) are all electrical or gas building loads that do not fall under lighting, heating, cooling, ventilation, or water heating. As lighting and HVAC equipment becomes more efficient, PPLs represent an increasing percentage of commercial building energy consumption---up to 50\% in high-performance buildings. This section describes how electric equipment, gas equipment, data centers, elevators, and kitchen equipment are modeled in ComStock. 

\subsection{Electric Equipment}

Electric equipment is a broad category that encompasses any type of load that is powered by an AC outlet. This can include computers, monitors, printers, kitchen or bathroom appliances, laundry equipment, phone chargers, and more. Because there are so many types of electric equipment, ComStock does not model each technology separately, but instead uses an equipment power density (EPD) value for each space type, in watts per square foot. 

The EPDs are a property of the space type. Each space type names an equipment category, and the equipment data carries one or more rows for that category. Where the source data distinguishes the building type the space is in---a classroom in a primary school carries less plug load than one in an office building, and a dining room in a hotel more than one in a school---the row matching the model's primary building type is used; where the data does not distinguish, a single row applies to every building type; and a space type with no row for its building type uses the median across the building types that have one. The values derive from the DOE prototype buildings, with the adjustments made during the end-use load profiles calibration process carried into the data. Using end-use-level data provided by two industry sources for a variety of building types, we increased or decreased some EPD assumptions to better reflect the actual building data. Building types that were affected by the EPD adjustments included Full Service Restaurant (FullServiceRestaurant), Primary School (PrimarySchool), Quick Service Restaurant (QuickServiceRestaurant), Retail (Retail), Secondary School (SecondarySchool), and Strip Mall (StripMall).

Equipment power density does not vary with the building's template. Plug and process loads describe how a space is used rather than what it was legal to install when the building was built, and the typical data is therefore vintage-agnostic; the interior equipment template continues to govern only the efficiency of the equipment that the standards data describes. The EPDs in the ComStock model for each space type, and the building type rows where they differ, are shown in Table~\ref{tab:epds}.


\subsection{Gas Equipment}

Gas equipment refers to any natural gas-powered interior equipment that is not used for space heating or water heating. Similar to electric equipment, there are many different types of gas equipment, so ComStock does not model each technology individually, but rather uses a gas intensity in BTU per hour per square foot, keyed on the space type in the same way as electric equipment. Gas kitchen equipment makes up the majority of the gas equipment modeled in ComStock. Kitchen equipment is discussed separately in Section~\ref{sec:kitchen_equipment}; the kitchen values in the data are placeholders that the kitchen equipment measure replaces with counted appliances. Outside kitchens, only the laundry and operating room space types carry gas equipment. The values used are shown in Table~\ref{tab:gas_equip}.

\begin{table}[h!]
\scriptsize
\centering
\caption[Gas Equipment Power Density]{Gas equipment power density (Btu/hr-ft\textsuperscript{2}) by space type. Kitchen (food preparation) values are the placeholder loads that the kitchen equipment measure replaces with counted appliances, as described in Section~\ref{sec:kitchen_equipment}.}
\label{tab:gas_equip}
\begin{tabular}{|p{2.2in}|p{1.3in}|r|r|r|r|}
\hline
\textbf{Space Type} & \textbf{Building Type} & \textbf{Btu/hr-ft\textsuperscript{2}} & \textbf{Radiant} & \textbf{Latent} & \textbf{Lost} \\ \hline
food preparation & FullServiceRestaurant & 205.8 & 0.2 & 0.1 & 0.7 \\ \hline
 & Hospital & 51.2 & 0.5 & 0 & 0 \\ \hline
 & LargeHotel & 767 & 0.2 & 0.1 & 0.7 \\ \hline
 & QuickServiceRestaurant & 250.9 & 0.2 & 0.1 & 0.7 \\ \hline
 & SuperMarket & 8.5 & 0.2 & 0.1 & 0.7 \\ \hline
food preparation - deli/bakery & any & 100 & 0.2 & 0.1 & 0.7 \\ \hline
food preparation - primary school & any & 453.7 & 0.2 & 0.1 & 0.7 \\ \hline
food preparation - secondary school & any & 532.1 & 0.2 & 0.1 & 0.7 \\ \hline
laundry/washing & LargeHotel & 170 & 0.5 & 0 & 0 \\ \hline
 & SmallHotel & 129.9 & 0.2 & 0 & 0 \\ \hline
laundry/washing - coin-operated laundries & LargeHotel & 170 & 0.5 & 0 & 0 \\ \hline
 & SmallHotel & 129.9 & 0.2 & 0 & 0 \\ \hline
operating room & any & 23.9 & 0 & 0 & 0.5 \\ \hline
\end{tabular}
\end{table}


Gas laundry equipment, in the laundry space types of large and small hotels, represents gas clothes dryers, which are common in commercial drying applications. In operating rooms, in hospitals and outpatient buildings, a small amount of gas equipment represents steam sterilizers or autoclaves, which are used for sterilization during surgical procedures.

\subsection{Data Centers}

Data centers are a high-intensity type of PPL that house IT and computing equipment. Large standalone data centers are not currently modeled in ComStock, but this building type may be added in the future. Instead, ComStock models data centers as a space type within large and medium office buildings. This is meant to represent an IT closet or high-performance computer that is located within an office building and used by a business or organization. 

Data centers are represented by two space types---a high-ITE data center and a low-ITE data center. The high-ITE space type represents the main computer room and has an equipment power density of 45 W/ft\textsuperscript{2}, approximately 20--50 times the equipment load of most other space types \citep{osti_1129366}; the low-ITE space type represents the IT closets and support space around it at a much lower density. The equipment power densities of both are listed in Table~\ref{tab:epds}, and the fraction of floor area each occupies in the office definitions is listed in Table~\ref{tab:space_type_ratios_all}.

In the DOE large office prototype model, the data center represents 2.5\% of the total floor area. The medium office prototype model does not contain a data center space type. If we used the exact space type ratios from the prototype models for ComStock, all large offices would contain a data center, but no medium or small offices would contain this space type. In reality, not all office buildings contain data centers, and they can be present in offices of different sizes. Therefore, ComStock models data centers in a portion of large and medium office buildings, through building subtypes with and without the data center space types. To determine these distributions, we used CBECS 2012 data to understand how prevalent data centers are in office buildings of different sizes. In addition, we used CBECS responses to determine what percent of a typical office building's floor area is dedicated to the data center. From this analysis, we decided that 38\% of large offices and 20\% of medium offices should contain data centers. We also determined that data centers are uncommon in small offices; therefore, there is no data center space type in the small office models.

Data centers follow a very different schedule than the rest of a building's plug and process loads. This type of IT equipment often runs 24 hours a day; therefore, the data center space types have no occupancy and a flat equipment schedule year round, and they are flagged as critical operation so that their zone HVAC equipment runs continuously. The hours of operation and base-to-peak ratio (BPR) adjustments do not affect the data center space types, and no sampled thermostat setpoint or setback is applied to them: they hold the data center setpoints in the space type data around the clock.

\subsection{Elevators}
Elevators are a high power density equipment load present in many commercial buildings. According to the Americans with Disabilities Act (ADA), elevators are required in all commercial buildings with three or more stories, or when the square footage of each floor is more than 3,000 square feet. Although not a requirement, many two-story buildings also contain elevators for accessibility and convenience. Therefore, ComStock includes elevators in all buildings with two or more stories.

Hydraulic elevators are assumed to be installed in buildings with two to six stories, and traction elevators are assumed to be installed in buildings taller than six stories. Hydraulic elevators use a fluid-driven piston to lift the cab, and typically operate at speeds of 150 feet per minute or less. Hydraulic elevators are more affordable and can carry heavier loads, but because of their slow speeds, they are typically only used in buildings up to five stories. For buildings with six or more stories, traction elevators are used because they operate at much higher speeds (up to 500 feet per minute). Traction elevators use a counterweight and pulley system, making them more energy efficient because the motor does not have to move as much weight. The drawbacks, however, include high installation and maintenance costs and limits on cab weights. The motor power is assumed to be 16,055 W for hydraulic elevators and 20,370 W for traction elevators. 

Elevators are modeled as a zone load in EnergyPlus, meaning the elevator equipment load and associated heat gain are attributed to a thermal zone. Elevator load is reported out as part of the electric equipment end use. With hydraulic elevators, the elevator room is typically located in the basement, so the equipment load and heat gain are added to the first floor core zone. With traction elevators, the elevator equipment is located on the roof, so the equipment load and heat gain are added to the top floor core zone.

The number of elevators installed in a ComStock building depends on the building type. For most building types, the number of elevators is based on the floor area. The exceptions are hospitals and hospitality buildings, for which assumptions are based on the number of hospital beds or hotel rooms, respectively. ComStock differentiates between passenger elevators and freight elevators in order to properly capture the elevator load in certain building types with industrial or service elevators. Passenger elevators are modeled in all building types, whereas freight elevators are only modeled in hospital, large hotel, large office, and warehouse buildings. The assumptions for the number of passenger and freight elevators modeled by building type are shown in Tables ~\ref{tab:passenger_elevators} and ~\ref{tab:freight_elevators}.

The equipment schedule for elevators is irregular and unpredictable. Therefore, this load does not follow the typical plug load schedule for its associated space type. We decided to approximate an elevator's schedule by relating it to the number of people who are entering or exiting the building at each time step---in other words, the derivative of the occupancy schedule of the building. We also made assumptions regarding the number of people per elevator ride (five), the amount of time per ride (calculated from the elevator speed and number of stories), and the amount of inter-floor traffic that is not captured by the change in building occupancy (added a factor of 1.2x). Elevator data and metrics like this are not commonly measured or available, so these assumptions are based primarily on engineering judgment. 

From these calculations and assumptions, we derived an elevator schedule for each building type and each day of the week. An example of the elevator schedule for medium offices is shown in Figure~\ref{fig:medium_office_elevator_schedule} for weekdays, Saturdays, and Sundays. On weekdays, the most elevator traffic occurs at the beginning and end of the day, when people are coming in or leaving work for the day. There is also significant traffic during the lunch hour, as some people choose to leave the building for lunch. At all other times during the workday, the elevator load is approximately 40\% of the total load to account for inter-floor traffic and minimal change in the total building occupancy. For some buildings, the occupancy schedules are reduced on Saturdays and include no occupancy on Sundays, which is reflected in the elevator schedule. 

\begin{figure}[ht!]
    \centering 
    \includegraphics[]{figures/medium_office_elevator.png}
    \caption[Medium office elevator schedule]{Elevator fractional load schedule for medium offices on weekdays, Saturdays, and Sundays.}
    \label{fig:medium_office_elevator_schedule}
\end{figure}

The final aspect of modeling elevators is accounting for lighting and fans inside the elevator. Although these are minimal loads compared to the total elevator energy, they are still modeled in ComStock. The elevator lighting and fans are defined in the model by a total power in watts. These wattage values were calculated based on the number of elevators and the subsystem template in the model, thereby ensuring higher efficiencies for newer lighting and ventilation systems. The elevator lighting and fan schedules are assumed to be at full load at all times when the elevator schedule is greater than 0.

\pagebreak

\subsection{Kitchen Equipment}
\label{sec:kitchen_equipment}
Kitchens are one of the most energy-intensive space types. In ComStock, the kitchen is the food preparation space type, which appears in six building types---FullServiceRestaurant, Hospital, LargeHotel, PrimarySchool, QuickServiceRestaurant, and SecondarySchool---and, as the deli and bakery food preparation space type, in grocery stores. In some building types, namely restaurants, the kitchen space type represents a significant proportion of the floor area. In these cases, kitchen loads have a major impact on the total building EUI. In hotels, hospitals, and schools, the kitchen space type only represents a small fraction of the total floor area. Table~\ref{tab:kitchen_floor_area} shows the percent of the total floor area represented by the kitchen space type for each building type. Strip malls that are sampled with a restaurant fraction (Section~\ref{sec:building_type_meta}) carry the full service restaurant's dining and food preparation space types in that fraction of their floor area, and receive cooking equipment in the same way as a restaurant. The cooking equipment described below is placed in every food preparation space in a model, whichever building type contains it, except the grocery deli and bakery service areas, which hold no primary cooking appliances and keep the gas and electric equipment densities of their space type.

\begin{table}[h!]
\small
\centering
\caption[Kitchen Space Type Percentage of Total Floor Area]{Kitchen (food preparation) space type percentage of total floor area}
\label{tab:kitchen_floor_area}
\begin{tabular}{|l|l|l|}
\hline
\textbf{Building Type} & \textbf{Space Type} & \textbf{Percentage of Total Floor Area} \\ \hline
Full Service Restaurant & food preparation & 27.3\% \\ \hline
Hospital & food preparation & 4.1\% \\ \hline
Large Hotel & food preparation & 0.9\% \\ \hline
Primary School & food preparation - primary school & 2.4\% \\ \hline
Quick Service Restaurant & food preparation & 50\% \\ \hline
Secondary School & food preparation - secondary school & 1.1\% \\ \hline
Grocery & food preparation - deli/bakery & 5\% \\ \hline
\end{tabular}
\end{table}


Commercial building energy modeling often assumes an energy intensity per area for cooking equipment, like what is used in the DOE prototype buildings (\cite{qsr_50pct_svs}). However, applying these power densities directly to the various building sizes in ComStock assumes cooking loads scale linearly with kitchen square footage, which does not necessarily represent reality (\cite{eia2012cbecs},~\cite{qsr_50pct_svs}). Further, it does not allow for straightforward analysis of specific cooking equipment since it is represented by a single aggregate load. However, little information exists regarding representative equipment-specific modeling of cooking appliances, as well as how cooking equipment scales with building area. The CBECS survey does provide some level of guidance for area-based intensity as they disaggregate cooking energy consumption separately and provide building area for their samples, but this end use energy consumption data is derived using statistical means from monthly billing data rather than directly measured making it less reliable (\cite{eia2012cbecs},~\cite{qsr_50pct_svs}).


ComStock uses published data to create representative probability distributions of commercial cooking equipment counts, by building type, for both gas and electric appliances. Additionally, the equipment distributions are scaled by area to represent the non-linear scaling suggested in the literature. Although other building types likely include some degree of cooking equipment as well, such as larger offices (\cite{eia2012cbecs}), the current implementation of ComStock only includes cooking equipment in the previously-mentioned six building types plus quick service restaurants found in strip malls. 

Commercial kitchens can contain electric or gas cooking equipment, or a mix of both. The prevalence of gas and electric fuel types for each equipment type used in ComStock are derived from a DOE study (\cite{goetzler_commercial_appliances}). ComStock requires rated input power values and fractions of radiant, latent, and lost heat for gas and electric kitchen equipment. These values are primarily derived from the ASHRAE Fundamentals Handbook (\cite{ashrae2017}) after comparisons with other kitchen equipment studies and commercially available products. More details about how these values were determined can be found in the End Use Savings Shapes documentation (\cite{nrel89130}).The assumptions used in ComStock for prevalence, rated input power, and fractions radiant, latent, and lost for gas and electric appliances are shown in Table~\ref{tab:kitchen_prev_and_power}. 

\input{tables/kitchen_prev_and_power}

Kitchens in ComStock models are assigned a quantity of each equipment type. These can be found in the ComStock metadata files for each model. Equipment quantities are assigned using probability distribtuions with dependencies on building type and floor area by food service and building type. The quantities used in the probability distributions are determined using an older dataset of equipment counts per restaurant type combined with the prevalence of the restaurant type (\cite{com_food_service_equip}). The restaurant types were mapped to ComStock building types. This is summarized in Tables~\ref{tab:kitchen_cook_counts}. Unfortunately, little data of this type was found in literature, so this older source was ultimately used for equipment counts. However, even with changes in culinary trends, it is not expected that counts of equipment types in a restaurant would have changed drastically over the past few decades. Additionally, a “None” restaurant type was created for schools, with a prevalence determined by the fraction of schools in CBECS that have kitchens (\cite{eia2012cbecs}).

This workflow also includes modifiers to the equipment counts shown in Table~\ref{tab:kitchen_cook_counts} to scale equipment count for different kitchen sizes. As mentioned previously, it is not expected that equipment scales linearly with kitchen size, but there is little information in the literature to suggest appropriate scaling (\cite{qsr_50pct_svs}). To account for equipment count scaling, we assume most typically-sized kitchens will have the same quantity of equipment. However, especially small and large kitchens will include scaling factors to account for very large changes in kitchen area that would likely correlate to higher/lower meals served. These factors were determined using engineering judgment and are summarized in Figure~\ref{fig:cooking_equipment_count_scaling_factors_by_building_type_and_area}. Additional scaling factors beyond those shown in Figure~\ref{fig:cooking_equipment_count_scaling_factors_by_building_type_and_area} were applied to Quick Service Restaurants and Full Service Restaurants below 10,000 square feet, as the initial scaling factors resulted in unrealistically high EUIs for very small restaurants. The following supplemental scaling factors were applied to equipment counts:
\begin{itemize}
\item Full Service Restaurants (all restaurant types): \\
    \begin{itemize}
    \item Under 5,000 square feet: equipment counts reduced by an additional 50\% \\
    \item Between 5,000 and 10,000 square feet: equipment counts reduced by an additional 30\% \\
    \end{itemize}
\item Quick Service Restaurants (all restaurant types): \\
    \begin{itemize}
    \item Under 5,000 square feet: equipment counts reduced by an additional 10\% \\
    \item Between 5,000 and 10,000 square feet: equipment counts reduced by an additional 50\%. \\
    \end{itemize}
\end{itemize}

In summary, ComStock determines quantity and fuel type of cooking equipment based on sampling our probability distributions. A restaurant type is sampled for each model as per the prevenances shown by building type in Table~\ref{tab:kitchen_cook_counts}. This yields the corresponding equipment counts for the restaurant type shown in Table~\ref{tab:kitchen_cook_counts}, with square footage modifiers applied based on building size and type. Next, the equipment fuel type is sampled as per the prevalence shown in Table~\ref{tab:kitchen_prev_and_power} for each piece of equipment. Combined, this yields quantities of gas and electric cooking equipment for each ComStock sample. Note that kitchen spaces in ComStock additionally include some prevalence of electric load to account for non-major electrical appliances such as microwaves, heating lamps, toasters, coffee machines, electric kettles, etc. Additionally, the current implementation of cooking equipment in ComStock utilizes the same schedule for each equipment type. Future work could include implementing equipment-specific schedules for the various equipment types, which may better represent reality.


\begin{figure}[ht!] \centering 
    \includegraphics[
        page={1},
        %trim={1cm 18.3cm 1cm 1cm}, clip, % L B R T
        width=\textwidth]{figures/kitchen_equipment_count_scaling.png}
    \caption[Cooking equipment count scaling factors by building type and area]{Cooking equipment count scaling factors by building type and area.}
    \label{fig:cooking_equipment_count_scaling_factors_by_building_type_and_area}
\end{figure}
